\section{Related Work} 
\label{sec:related}

Our framework, along with its underlying formalism, analysis approach, and application domain, is closely related to the following three lines of  earlier research.

\paragraph{Formal Modeling and Analysis of DNS}
Kakarla et al.~\cite{groot} formalize the DNS resolution semantics, abstracting away critical resolver-side logic including caching and recursive queries. 
Building upon this semantics, they also 
provide the first static analyzer GRoot for identifying DNS misconfiguration errors. 
Liu et al.~\cite{dnsmaude} complement GRoot's semantics by
providing the first formal 
semantics of end-to-end name resolution.
The accompanying tool DNSMaude
offers a collection of formal analyses for both qualitative (e.g., rewrite blackholing)
and quantitative 
properties (e.g., amplification).
We have already discussed
their limitations in \Cref{sec:intro}. 




\paragraph{Testing, Monitoring, and Verification of DNS Implementations} 
Recent years have seen significant advances in DNS testing tools~\cite{Eywa,ResolFuzz,ResolverFuzz,scale}. 
ResolFuzz~\cite{ResolFuzz} is a differential tester, searching for semantic
bugs such as inaccuracies in resolver responses across DNS resolver implementations.
ResolverFuzz~\cite{ResolverFuzz} fuzzes resolvers via
 short query-response sequences. 
 It also leverages differential testing and can identify resolver vulnerabilities like cache poisoning bugs.
Unlike pure fuzzing,
SCALE~\cite{scale}
guides its test generation by  
symbolically executing GRoot's  semantics. 
It focuses on detecting RFC compliance errors in authoritative nameserver implementations rather than recursive resolvers.
While also deriving exhaustive test cases using symbolic execution, Eywa~\cite{Eywa} leverages large language models to construct intended DNS behaviors from natural language sources like RFCs. 
All these testers, along with excellent DNS monitors~\cite{thousandeyes,check-host}, 
can only show the presence of bugs or vulnerabilities, not their absence, in DNS implementations.
Moreover, in contrast to our framework, they
do not
encompass the full name resolution process involving both  recursive resolvers and nameservers.


IRONSIDES~\cite{Ironsides} is a nameserver
implementation provably secure against DNS vulnerabilities like single-packet DoS
attacks. 
However, it may still be susceptible to application-layer issues such as amplification attacks, in contrast to our framework. 
DNS-V~\cite{DNS-V} is a verification
framework for an in-house DNS authoritative 
engine in Alibaba Cloud.
A symbolic-execution-based approach is used to 
facilitate the adoption of layered verification, where 
source code within each decomposed layer of DNS can be verified independently with respect to correctness properties. 
Quantitative properties like amplification are out of its scope.
Wang et al.~\cite{10.1145/3663408.3663412}  scale up the verification of DNS configuration
by parallelizing the analysis of zone-file-specific query processing behaviour and symbolically combining the analysis results. 
Unlike our framework, the use of formal methods in none of these works is related to the DNS semantics.


All these testing, monitoring, and verification efforts focus on DNS implementations.
Given that many DNS vulnerabilities have stemmed
 from RFCs or protocol-level defects in heavily tested production resolvers and nameservers, protocol-level verification
 at an early design stage is highly desirable. Our framework takes a step forward towards this goal. 



\paragraph{Verification Problem for Infinite-State Systems}
Verification of infinite-state systems is a well-studied problem and is still an active area of research. Abdulla \& Jonsson \cite{AbJo:lossy,AbJo:lossy:IC} and C{\'e}c{\'e} et al. \cite{cece1996unreliable} independently show that the verification problem for CPSes communicating over lossy channels is decidable in non-primitive recursive time, for arbitrary topologies. Abdulla et al. \cite{abdulla2012timed} later show that the decidability holds even when the communicating processes are timed and the messages contain time-stamps. Heussner et al. \cite{heussner2010reachability,heussner2012reachability} study this problem for rCPSes, however, over perfect channels. They restrict their attention to eager runs and show that the problem is decidable in exponential-time for a restricted class of topologies that they call \emph{non-confluent}. Prior to this work, Atig et al. \cite{atig2008reachability} show that the problem is decidable even without the restriction of eagerness, but instead requires the underlying topology to be acyclic. In the context of PDSes, Bouajjani et al. \cite{bouajjani1997reachability} show that the problem is decidable for PDS, using a \textit{saturation method} that computes the set of all reachable computations. Abdulla et al. \cite{abdulla2012dense} and Clemente and Lasota \cite{clemente2015timed} extend the decidability to timed PDS with a timed stack. In contrast, our work handles a timed extension of rCPS with an infinite message alphabet. In addition, we focus on models pertaining to the DNS semantics.

Kakarla et al. \cite{dns-complex} study the expressive complexity of DNS. They show that the verification of DNS zone files, a simpler problem than $\verif$, is likely to take at least cubic time in
the number of records, effectively giving a lower bound. In this work, we provide an upper bound for this problem (\Cref{lem:complexity}).






    

























